\documentclass[11pt,a4paper]{ctexart}
\usepackage[utf8]{inputenc}
\usepackage{amsmath}
\usepackage{amssymb}
\usepackage{geometry}
\geometry{left=2.5cm,right=2.5cm,top=2.5cm,bottom=2.5cm}

\title{\textbf{二维双曲 $CL$ 网格阻抗积分的留数定理详细推导}}
\author{物理与电路理论推导}
\date{\today}

\begin{document}

\maketitle

\section{问题背景与基尔霍夫方程定义}
考虑一个二维无穷网格电路，其中横向（$x$ 方向）全部由电容 $C$ 连接，纵向（$y$ 方向）全部由电感 $L$ 连接。
设在节点 $(0,0)$ 注入交流电流 $I$，在节点 $(m,n)$ 流出电流 $-I$。任意节点 $(x,y)$ 的电势记为 $V(x,y)$。

根据基尔霍夫电流定律（KCL），对任意节点 $(x,y)$ 列出流出电流方程：
\begin{equation}
\frac{V(x,y) - V(x-1,y)}{\frac{1}{i\omega C}} + \frac{V(x,y) - V(x+1,y)}{\frac{1}{i\omega C}} + \frac{V(x,y) - V(x,y-1)}{i\omega L} + \frac{V(x,y) - V(x,y+1)}{i\omega L} = I(x,y)
\end{equation}
其中源电流分布定义为：
\begin{equation}
I(x,y) = \begin{cases} 
I, & x=0, y=0 \\
-I, & x=m, y=n \\
0, & \text{其他}
\end{cases}
\end{equation}

将等式 (1) 两边同乘以 $-i$，整理算子项得到：
\begin{equation}
\omega C \big[ 2V(x,y) - V(x-1,y) - V(x+1,y) \big] - \frac{1}{\omega L} \big[ 2V(x,y) - V(x,y-1) - V(x,y+1) \big] = -i I(x,y)
\end{equation}

\section{离散傅里叶变换与频域求解}
引入二维离散傅里叶变换（Discrete Fourier Transform），定义空间域到动量域的变换对为：
\begin{equation}
V(k_1, k_2) = \sum_{x=-\infty}^{+\infty} \sum_{y=-\infty}^{+\infty} V(x,y) e^{i(x k_1 + y k_2)}
\end{equation}
对源电流 $I(x,y)$ 进行变换，由于只有原点和 $(m,n)$ 点有值，得到：
\begin{equation}
\tilde{I}(k_1, k_2) = I \big( 1 - e^{i(m k_1 + n k_2)} \big)
\end{equation}

将离散拉普拉斯算子变换到差分频域：
\begin{align}
2V(x,y) - V(x-1,y) - V(x+1,y) &\xrightarrow{\mathcal{F}} 2(1 - \cos k_1) V(k_1, k_2) = 4\sin^2\frac{k_1}{2} V(k_1, k_2) \\
2V(x,y) - V(x,y-1) - V(x,y+1) &\xrightarrow{\mathcal{F}} 2(1 - \cos k_2) V(k_1, k_2) = 4\sin^2\frac{k_2}{2} V(k_1, k_2)
\end{align}

将变换后的算子代入总方程 (3) 中：
\begin{equation}
\left[ 4\omega C \sin^2\frac{k_1}{2} - \frac{4}{\omega L} \sin^2\frac{k_2}{2} \right] V(k_1, k_2) = -i I \big( 1 - e^{i(m k_1 + n k_2)} \big)
\end{equation}
两边同乘 $\omega L$ 并移项，解出频域电压响应 $V(k_1, k_2)$：
\begin{equation}
V(k_1, k_2) = \frac{\omega L I \big( 1 - e^{i(m k_1 + n k_2)} \big)}{4i \left( \omega^2 L C \sin^2\frac{k_1}{2} - \sin^2\frac{k_2}{2} \right)}
\end{equation}

\section{目标阻抗积分的建立}
任意节点 $(m,n)$ 与原点 $(0,0)$ 之间的等效阻抗 $\tilde{Z}(\omega)$ 定义为：
\begin{equation}
\tilde{Z}(\omega) = \frac{V(m,n) - V(0,0)}{I}
\end{equation}
利用傅里叶逆变换 $V(x,y) = \frac{1}{4\pi^2}\int_{-\pi}^{\pi}\int_{-\pi}^{\pi} V(k_1,k_2) e^{-i(xk_1+yk_2)} \,dk_1 dk_2$，我们可以将阻抗写为双重格林函数积分形式：
\begin{equation}
\tilde{Z}(\omega) = \frac{-i\omega L}{16\pi^2} \int_{-\pi}^{\pi} \int_{-\pi}^{\pi} \frac{\big(1 - e^{i(m k_1 + n k_2)}\big)\big(e^{-i(m k_1 + n k_2)} - 1\big)}{\omega^2 L C \sin^2\frac{k_1}{2} - \sin^2\frac{k_2}{2}} \, dk_1 dk_2
\end{equation}
通过晶格对称性化简，或者直接考察单端注入电势响应差，核心的非零贡献积分项可写为：
\begin{equation}
\tilde{Z}(\omega) = \frac{-i\omega L}{4\pi^2} \int_{-\pi}^{\pi} \int_{-\pi}^{\pi} \frac{1 - e^{i(m k_1 + n k_2)}}{\omega^2 L C \sin^2\frac{k_1}{2} - \sin^2\frac{k_2}{2}} \, dk_1 dk_2
\end{equation}

\section{利用留数定理进行解析降维化简}
为了解决分母由于减号在实数轴形成的双曲共振奇点线，利用半角公式展开分母：
\begin{equation}
\omega^2 L C \sin^2\frac{k_1}{2} - \sin^2\frac{k_2}{2} = \frac{\omega^2 L C (1 - \cos k_1) - (1 - \cos k_2)}{2}
\end{equation}
代入阻抗公式 (12) 并消去常数系数，将积分写为关于 $k_2$ 的嵌套形式：
\begin{equation}
\tilde{Z}(\omega) = \frac{-i\omega L}{2\pi^2} \int_{-\pi}^{\pi} \left[ \int_{-\pi}^{\pi} \frac{1 - e^{i m k_1} e^{i n k_2}}{\omega^2 L C (1 - \cos k_1) - 1 + \cos k_2} \, dk_2 \right] dk_1
\end{equation}
定义只与 $k_1$ 相关的独立参量：
\begin{equation}
A(k_1) = \omega^2 L C (1 - \cos k_1) - 1
\end{equation}
记内层关于 $k_2$ 的积分为 $J(k_1)$：
\begin{equation}
J(k_1) = \int_{-\pi}^{\pi} \frac{1 - e^{i m k_1} e^{i n k_2}}{A(k_1) + \cos k_2} \, dk_2
\end{equation}

\subsection{复平面围道变换}
令复变量 $z = e^{i k_2}$，当 $k_2 \in [-\pi, \pi]$ 时，积分路径转化为复平面上的单位圆周 $C$（逆时针方向）。此时有：
\begin{equation}
e^{i n k_2} = z^n, \quad \cos k_2 = \frac{z + z^{-1}}{2}, \quad dk_2 = \frac{dz}{iz}
\end{equation}
代入 $J(k_1)$ 中，整理为标准的有理分式复围道积分：
\begin{equation}
J(k_1) = \oint_{C} \frac{1 - e^{i m k_1} z^n}{A(k_1) + \frac{z + z^{-1}}{2}} \frac{dz}{iz} = \oint_{C} \frac{2\big(1 - e^{i m k_1} z^n\big)}{i\big(z^2 + 2A(k_1)z + 1\big)} \, dz
\end{equation}

\subsection{极点分析与留数计算}
令分母二次方程 $z^2 + 2A(k_1)z + 1 = 0$ 为零，求得两个孤立极点：
\begin{equation}
z_\pm = -A(k_1) \pm \sqrt{A(k_1)^2 - 1}
\end{equation}
因为 $z_+ \cdot z_- = 1$，在引入微小因果损耗 $\omega \to \omega + i\epsilon$ 后，必有一个极点落在单位圆内部（$|z| < 1$），另一个落在圆外。我们选取圆内的有效拓扑极点记为 $z_0$：
\begin{equation}
z_0 = -A(k_1) + \text{sgn}\big(A(k_1)\big)\sqrt{A(k_1)^2 - 1}
\end{equation}
根据留数定理，该一阶极点处的留数（Residue）计算如下：
\begin{align}
\text{Res}\left( \frac{2(1 - e^{i m k_1} z^n)}{i(z^2 + 2A(k_1)z + 1)}, z_0 \right) &= \left. \frac{2(1 - e^{i m k_1} z^n)}{i \frac{d}{dz}(z^2 + 2A(k_1)z + 1)} \right|_{z=z_0} \nonumber \\
&= \frac{2(1 - e^{i m k_1} z_0^n)}{i(2z_0 + 2A(k_1))} \nonumber \\
&= \frac{1 - e^{i m k_1} z_0^n}{i \big(z_0 + A(k_1)\big)}
\end{align}
将 $z_0 + A(k_1) = \text{sgn}\big(A(k_1)\big)\sqrt{A(k_1)^2 - 1}$ 代入，得：
\begin{equation}
\text{Res} = \frac{1 - e^{i m k_1} z_0^n}{i \cdot \text{sgn}\big(A(k_1)\big)\sqrt{A(k_1)^2 - 1}}
\end{equation}

根据留数定理公式 $J(k_1) = 2\pi i \cdot \text{Res}$，虚数单位 $i$ 被精准消去：
\begin{equation}
J(k_1) = \frac{2\pi \cdot \text{sgn}\big(A(k_1)\big) \big(1 - e^{i m k_1} z_0^n\big)}{\sqrt{A(k_1)^2 - 1}}
\end{equation}

\section{最终的一维解析化简结论}
将求导完内圈积分的 $J(k_1)$ 代回总阻抗方程 (14) 中，消去常数系数，最终得到降维后的一维精确解析积分表达式：
\begin{equation}
\tilde{Z}(\omega) = \frac{-i\omega L}{2\pi} \int_{-\pi}^{\pi} \frac{\text{sgn}\big(A(k_1)\big) \left[ 1 - e^{i m k_1} \left( -A(k_1) + \text{sgn}\big(A(k_1)\big)\sqrt{A(k_1)^2 - 1} \right)^n \right]}{\sqrt{A(k_1)^2 - 1}} \, dk_1
\end{equation}
其中核心依赖参量为：
\begin{equation}
A(k_1) = \omega^2 L C (1 - \cos k_1) - 1
\end{equation}
该公式成功消除了原级数展开带来的组合数发散困难，将原双重积分压缩为单重规则积分，极大地优化了数值计算的收敛速度和稳定性。






\end{document}